% Finalizado 0 
\chapter{Metodología} 
\label{sec:metodologia}

Para alcanzar el objetivo, es necesario adoptar un sistema en paralelo que aproveche las capacidades de la GPU. A continuación, se exponen las bases teóricas que fundamentan el desarrollo del motor.


\section{Definición de Partícula}
Una partícula es la unidad básica que se utiliza para simular fluidos, gases o cualquier fenómeno que carezca de una forma definida.
Dado que carecen de una malla de triángulos compleja, es suficiente con almacenar un número reducido de variables esenciales. Estas propiedades cambian en cada fotograma:
\begin{itemize}
    \item \textbf{Posición:} Tres coordenadas espaciales ($X, Y, Z$) para situar el elemento en el mundo 3D.
    \item \textbf{Velocidad:} Un vector dinámico. Indica la dirección y la fuerza del desplazamiento del elemento. Elementos como el viento o la gravedad modifican este vector constantemente.
    \item \textbf{Color y transparencia:} Controlan el aspecto visual. El motor interpola estos valores en función del tiempo. Esto permite observar cómo el humo se disipa o el fuego pierde intensidad.
    \item \textbf{Tiempo de vida (TTL):} Es la duración de la partícula medida en segundos. Cada elemento se genera en un emisor y finaliza su ciclo al consumir dicho tiempo, momento en el cual se reinicia. Esta técnica permite mantener un ciclo infinito sin necesidad de solicitar nueva asignación de memoria al sistema operativo en cada iteración.
\end{itemize}



\section{Sistemas de Partículas en CPU vs GPU}
Hay un gran reto para simular partículas: poder realizar millones de operaciones matemáticas por segundo.

El \textbf{enfoque clásico en CPU} trabaja de forma secuencial. El procesador calcula el estado de la partícula uno a uno. 
Al terminar, copia todo ese bloque de datos hacia la memoria de vídeo (VRAM) usando el bus PCI-Express. 
Este viaje constante satura el canal de comunicación. Se genera un cuello de botella crítico. 
Por este motivo, el rendimiento del procesador principal decae abruptamente si se intenta actualizar nubes masivas de elementos manteniendo la fluidez~\parencite{latta2004massively}.

Por el contrario, el enfoque basado en GPU (GPGPU) actualiza todo en paralelo. Básicamente aloja la información en la memoria gráfica 
y ejecuta la física en los \textit{Compute Shaders}~\parencite{akenine2018real}.

Como los datos ya no viajan por el bus PCI-Express, el rendimiento se dispara. La tasa de fotogramas sube de forma drástica.


\section{Fundamentos de Computación en GPU}


\subsection{El Compute Shader}
\label{subsection:ElComputeShader}

Esta etapa del pipeline es totalmente programable, lo que permite ejecutar algoritmos generales directamente en la GPU.
Esto sin interferir en el renderizado normal de otros elementos de la escena. 
Su fuerte es el cálculo matemático bruto. En lugar de procesar tareas pesadas en cadena, la tarjeta gráfica arranca miles de hilos ligeros a la vez. 
Constituye una herramienta idónea para el motor desarrollado. Cada hilo se encarga de evaluar y actualizar una única partícula de forma concurrente e independiente.


\subsection{El Pipeline de Computación}
El pipeline gráfico se encarga de transformar geometría en píxeles. 
Esto difiere del pipeline de computación, ya que este es un flujo lógico más simple que solo procesa datos.
Destaca por ser muy independiente. La GPU modifica millones de posiciones sin pedir ayuda a las etapas de dibujado. El proceso sigue tres pasos:
\begin{enumerate}
    \item \textbf{Despacho (\textit{Dispatch}):} Se define la cantidad de grupos de hilos necesarios para cubrir la totalidad de las partículas.
    \item \textbf{Cómputo en paralelo:} Los hilos ejecutan sus rutinas organizados en bloques, aprovechando de forma óptima la memoria compartida local. 
    \item \textbf{Sincronización:} La representación de los resultados en pantalla aguarda a la finalización de los cálculos, lo cual se garantiza mediante el uso de barreras de sincronización.
\end{enumerate}


\subsection{Técnica de Renderizado}
Para renderizar y computar simultáneamente sin merma de rendimiento, se emplea una estrategia de doble búfer en la GPU. Esto se conoce como \textit{Ping-Pong Buffers}~\parencite{akenine2018real}. 
Para ello se requieren dos bloques de memoria (véase la Tabla \ref{tab:ping_pong_buffers}). Este esquema simétrico agiliza el flujo interno y evita 
que un hilo intente leer y escribir en el mismo sitio.

\begin{table}[ht]
    \centering
    \renewcommand{\arraystretch}{1.2}
    \begin{tabular}{|l|p{0.7\textwidth}|}
        \hline
        \multicolumn{2}{|c|}{\textbf{Ping-Pong Buffers}} \\ \hline
        \textbf{Bloque de memoria} & \textbf{Descripción} \\ \hline
        \texttt{particlesIn} & Búfer de entrada que contiene el estado de las partículas del fotograma anterior (posiciones, velocidades y tiempo de vida restante) del cual lee el \textit{Compute Shader}. \\ \hline
        \texttt{particlesOut} & Búfer de salida donde el \textit{Compute Shader} escribe el nuevo estado de las partículas tras calcular las físicas. \\ \hline
    \end{tabular}
    \caption{Bloques de memoria en la estrategia de doble búfer (\textit{Ping-Pong Buffers}).}
    \label{tab:ping_pong_buffers}
\end{table} 

El funcionamiento de esta técnica se estructura en las siguientes fases:
\begin{enumerate}
    \item \textbf{Lectura de estado (\textit{Read}):} El \textit{Compute Shader} accede al búfer de entrada (\texttt{particlesIn}) para recuperar el estado anterior del sistema (posiciones, velocidades y tiempo de vida).
    \item \textbf{Procesamiento físico (\textit{Compute}):} Se resuelven las ecuaciones físicas que determinan la evolución dinámica de las partículas en el fotograma actual.
    \item \textbf{Escritura de estado (\textit{Write}):} Los nuevos valores calculados se almacenan de forma concurrente en el búfer de salida (\texttt{particlesOut}).
    \item \textbf{Intercambio (\textit{Swap}):} Una vez que las barreras de sincronización garantizan la finalización de los cálculos en la GPU, los roles de los búferes se invierten (el búfer de salida pasa a ser el de entrada para el fotograma subsiguiente).
\end{enumerate}

Finalizado este ciclo lógico, la etapa gráfica accede al búfer que contiene el estado actualizado, extrayendo las posiciones para proceder al renderizado de la escena. Este mecanismo iterativo asegura una separación rigurosa y eficiente entre el cálculo físico paralelo y el dibujado gráfico continuo.


\subsection{Simulaciones Procedimentales}

Recrear fluidos reales con las ecuaciones matemáticas de Navier-Stokes en tiempo real es inviable~\parencite{akenine2018real}. 
El coste en cómputo es demasiado alto. Por ello, en gráficos por ordenador se opta preferentemente por simulaciones procedimentales. 
Estas técnicas no persiguen un realismo físico molecular; en su lugar, imitan el aspecto visual mediante fórmulas eficientes que generan caos y aleatoriedad controlada.

\noindent En programación GPGPU existen dos formas de generar este movimiento:

\begin{enumerate}
    \item \textbf{Ruido espacial estructurado (como Perlin o texturas de ruido):} Interpola gradientes pseudoaleatorios en una rejilla. 
    Da transiciones muy suaves para gases, pero el coste matemático es elevado. 
    Además, almacenar el ruido en texturas 3D saturaría el ancho de banda de la VRAM debido a las lecturas constantes que efectuaría cada hilo.
    \item \textbf{Funciones de dispersión tipo Hash:} Se utilizan operaciones puramente matemáticas a nivel de bit y transformaciones lineales eficientes. Estas funciones reciben un valor de entrada y retornan un número decimal normalizado.
\end{enumerate}

Los generadores \textit{hash} integrados en Compute Shaders no tienen rival para mover millones de elementos. 
La semilla es el identificador de cada hilo (\textit{Thread ID}). 
Con esto, cada unidad de cálculo genera sus propios desfases de manera autónoma.

Mediante la aplicación de pequeños valores de tiempo (\textit{offsets}), se generan turbulencias asíncronas en las tres dimensiones, rompiendo así cualquier simetría visual.
El cálculo íntegro se lleva a cabo en las unidades aritméticas (\textit{ALUs}) de la GPU, garantizando un consumo mínimo de recursos. 
Esto posibilita mantener una alta tasa de fotogramas mientras se simulan nubes de puntos con un comportamiento orgánico y fluido.

